\documentclass[11pt]{article}
\usepackage{amsmath,amssymb,amsthm,geometry,hyperref}
\geometry{margin=1in}
\newtheorem{theorem}{Theorem}
\setlength{\emergencystretch}{2em}
\title{A positive-power counterexample for Taylor intersection roots\\
\large A solution of MathDB 328557}
\author{Research notes in Principia Math In Progress}
\date{11 September 2026}
\begin{document}
\maketitle

\section{Task 1: the problem and its resolution}
For $f(z)=z^n$ and a positive integer $r<n$, let
\[
 P_c(z)=\sum_{k=0}^r\frac{f^{(k)}(c)}{k!}(z-c)^k
\]
be the order-$r$ Taylor polynomial at $c$. Conjecture 2 in Section 5.2 of
Horwitz's paper \cite{horwitz} asserts that, when $0<a<b$, every nonreal
root of $P_b-P_a$ has real part strictly between $a$ and $b$.
The following refutes that assertion for a positive integer power.

\begin{theorem}\label{counterexample}
Take $r=5$, $n=24$, $a=1$, and $b=2$. Then $P_2-P_1$ has a conjugate pair
of nonreal roots $z_+,z_-$ satisfying
\[
 2.00042<\operatorname{Re}z_\pm<2.00044,
 \qquad 0.29417<|\operatorname{Im}z_\pm|<0.29419.
\]
In particular, their real parts are greater than $b$.
\end{theorem}
All source hypotheses hold: $0<a<b$, $r>3$, $n\geq r+1$, and
$f^{(r+1)}$ is continuous and strictly positive on $[a,b]$.
The source's examples for negative powers concern a different extension.
The counterexample here uses the positive power explicitly required by
Conjecture 2. It does not refute the weaker Conjecture 1, which asks only
for at least one conjugate pair with real part between the nodes.

\section{The exact polynomial}
Put $y=z-2$. Directly from the binomial coefficients,
\begin{align*}
 P_2(2+y)&=\sum_{k=0}^5\binom{24}{k}2^{24-k}y^k,\\
 P_1(2+y)&=\sum_{k=0}^5\binom{24}{k}(1+y)^k.
\end{align*}
Their difference is $7Q(y)$, where
\begin{align*}
 Q(y)={}&3183470664y^5+1591706490y^4+606309440y^3\\
        &+165304680y^2+28723560y+2388823.
\end{align*}
Thus it suffices to isolate a nonreal zero of $Q$ in the right half-plane.

\section{An exact rational root certificate}
Let
\[
 c=\frac{43+29418i}{100000},\qquad \rho=\frac1{100000},\qquad
 d_k=\frac{Q^{(k)}(c)}{k!}.
\]
Every $d_k$ has rational real and imaginary parts. For a complex number $v$
write $U(v)=|\operatorname{Re}v|+|\operatorname{Im}v|$ and
$D(v)=\max(|\operatorname{Re}v|,|\operatorname{Im}v|)$; then
$D(v)\leq|v|\leq U(v)$. Exact substitution gives the following rational
inequalities, with deliberately loose integer bounds:
\[
\begin{array}{c|r}
 k&\text{upper bound for }U(d_k)\\\hline
 0&89\\
 2&937000000\\
 3&4028000000\\
 4&6282000000\\
 5&3184000000
\end{array}
\qquad D(d_1)>65000000.
\]
These are finite integer-arithmetic claims: the companion verifier evaluates
each rational real and imaginary part exactly and checks the displayed
inequalities. They can also be checked directly from the displayed $Q$ and $c$.

Since every upper bound in rows $2$ through $5$ is less than $10^{10}$,
\[
 \sum_{k=2}^5 |d_k|\rho^k
 <10^{10}\frac{\rho^2}{1-\rho}<2.
\]
On the circle $|y-c|=\rho$ it follows that
\[
 |Q(y)-d_1(y-c)|
 \leq |d_0|+\sum_{k=2}^5|d_k|\rho^k
 <91<650<|d_1(y-c)|.
\]
Rouch\'e's theorem therefore gives exactly one zero of $Q$ in
$|y-c|<\rho$, counted with multiplicity, because the comparison polynomial
$d_1(y-c)$ has one. The disk lies entirely in
\[
 0.00042<\operatorname{Re}y<0.00044,
 \qquad 0.29417<\operatorname{Im}y<0.29419.
\]
The coefficients of $Q$ are real, so its conjugate disk contains the
conjugate root. Adding $2$ gives the assertions of Theorem~\ref{counterexample}.
This closes the proof; decimal root approximations are not used.

\section{Follow-up tasks 2--5}
\begin{enumerate}
\setcounter{enumi}{1}
\renewcommand{\labelenumi}{\textbf{Task \arabic{enumi}.}}
\item Determine all pairs $(r,n)$ for which every nonreal intersection root
lies in $a<\operatorname{Re}z<b$ for every $0<a<b$. In particular, find the
least positive exponent producing failure for each fixed Taylor order.
\item For fixed $(r,n)$, classify the ratios $b/a>1$ at which a root crosses
either boundary of the vertical strip. Obtain exact algebraic criteria and
locate every transition for the case $(r,n)=(5,24)$.
\item Prove a precise limit theorem for the roots as $n\to\infty$ with $r$
and $b/a$ fixed. Determine how the zeros of truncated exponential series
control the number and size of excursions beyond $b$.
\item Resolve the weaker assertion that at least one nonreal conjugate pair
has real part in $(a,b)$ whenever $f^{(r+1)}$ is continuous and nonvanishing
there and $r\geq4$. Identify sharp hypotheses if this assertion also fails.
\end{enumerate}

\begin{thebibliography}{9}
\bibitem{horwitz} A. Horwitz, \emph{Means and nonreal intersection points of
Taylor polynomials}, Journal of Classical Analysis 6 (2015), 1--28,
Conjecture 2 in Section 5.2, p.~26.
\url{https://doi.org/10.7153/jca-06-01}.
Preprint: \url{https://arxiv.org/abs/1408.2573}.
\bibitem{mathdb} MathDB 328557, live record retrieved 11 September 2026.
\url{https://mathdb.com/p/328557}.
\end{thebibliography}
\end{document}
